\section{Порівняння криптографічних бібліотек}

\subsection{Обрані криптографічні бібліотеки}

Для проведення порівняльного аналізу було обрано дві найпопулярніші криптографічні бібліотеки мови програмування Python:
\begin{itemize}
    \item \textbf{cryptography} \cite{cryptography, cryptography git} --- офіційно рекомендована спільнотою Python Cryptographic Authority (PyCA) бібліотека. Вона побудована як надійна обгортка над перевіреними низькорівневими C-бібліотеками (насамперед OpenSSL) з безпечним внутрішнім ядром на мові Rust. Її архітектурна концепція чітко розмежовує простий високорівневий безпечний шар (безпечні рецепти типу \texttt{Fernet}) та низькорівневий шар примітивів \texttt{hazmat} (Hazardous Materials).
    \item \textbf{PyCryptodome} \cite{PyCryptodome, PyCryptodome git} --- самостійний форк занедбаної бібліотеки \texttt{PyCrypto}, створений для забезпечення зворотної сумісності, суттєвого розширення підтримуваних алгоритмів та виправлення критичних проблем безпеки й швидкодії. Вона є самодостатньою бібліотекою, що містить власні реалізації алгоритмів на мові C без жорсткої прив'язки до зовнішнього системного OpenSSL.
\end{itemize}

\subsection{Порівняльний аналіз бібліотек}

\begin{longtable}{|p{2cm}|p{6cm}|p{6cm}|}
\hline
\textbf{Критерій} &
\textbf{cryptography \cite{cryptography} and \cite{cryptography git}} &
\textbf{PyCryptodome \cite{PyCryptodome}, \cite{PyCryptodome git} and \cite{NVD PyCryptodome}} \\
\hline

Мови та платформи &
Мова: Python 3.9+, PyPy 3.11
Платформи: CentOS Stream 9–10 (x86-64); Fedora (x86-64); macOS 26 Tahoe (ARM64); 
Ubuntu 22.04, 24.04, 26.04+ (x86-64, ARM64, ARMv7l, ppc64le); Debian 12 Bookworm, 13 Trixie, Sid (x86-64); 
Alpine (x86-64, ARM64); Windows Server 2025 (64-bit).&
Мова: Python 3.8+, PyPy. Критичні операції реалізовані на C.
Платформи: Linux, Windows, macOS; архітектури x86, x86-64, ARM (включаючи ARM64 / Apple Silicon). Поширюється у вигляді попередньо скомпільованих бінарних коліс (wheels). \\
\hline

Симетрич\-не шифрування &
Підтримується. Бібліотека містить реалізації симетричних алгоритмів, зокрема AES, ChaCha20, 
Camellia, 3DES та ін., а також різні режими роботи блочних шифрів.&
Підтримується. Містить як сучасний AES (128, 192, 256 біт), Salsa20, ChaCha20, так і велику кількість застарілих/легасі-шифрів: DES, 3DES, Blowfish, CAST-128, RC2, ARC4. Підтримує режими ECB, CBC, CFB, OFB, CTR, OpenPGP. \\
\hline

AEAD &
Підтримується. Є готові реалізації AES-GCM, ChaCha20-Poly1305, AES-CCM, AES-SIV, AES-GCM-SIV та AES-OCB3.&
Підтримується. Реалізовано режими блокового автентифікованого шифрування GCM, CCM, EAX, SIV, OCB, а також комбінацію ChaCha20-Poly1305. Дозволяє передавати відкриті асоційовані дані (AAD). \\
\hline

Асимет\-рична криптографія &
Підтримується. Бібліотека містить RSA, Diffie–Hellman, ECDH, ECDSA, Ed25519, Ed448, X25519, X448,
а в актуальних версіях також постквантові ML-KEM та ML-DSA.&
Підтримується. Реалізовано алгоритми RSA (розмір ключів 1024--4096+ біт, схеми PKCS\#1 v1.5 та OAEP), DSA, а також еліптичні криві (ECC): стандартні криві NIST (P-256, P-384, P-521), Ed25519, X25519. Постквантові алгоритми відсутні. \\
\hline

Цифровий підпис &
Підтримуються RSA-PSS, RSA-PKCS1v1.5, ECDSA, Ed25519, Ed448 та ML-DSA.&
Підтримуються схеми RSA-PSS, RSA-PKCS1v1.5, DSA, ECDSA (згідно зі стандартом FIPS 186-4) та EdDSA (Ed25519). \\
\hline

Геш-функції та KDF &
Підтримуються  SHA-2, SHA-3, BLAKE2, SHA-1, MD5, SM3, а також XOF.&
Підтримуються SHA-1, SHA-2, SHA-3 (включно з SHAKE128/256, Keccak), BLAKE2 (2b, 2s), MD5; MAC-функції: HMAC, CMAC, Poly1305; KDF: scrypt, PBKDF2, HKDF, bcrypt. \\
\hline

Ініціаліза\-ція генератора ви\-падкових чисел &
Не потребує ручного налаштування PRNG користувачем. Випадкові значення генеруються 
за допомогою захищених механізмів операційної системи та OpenSSL. 
Тому користувачеві не потрібно самостійно додавати <<шум>> для ініціалізації генератора.&
Не потребує ручної ініціалізації та підмішування ентропійного <<шуму>> користувачем у просторі користувача. Модуль повністю покладається на системне джерело ентропії ОС. \\
\hline

Крипто\-графічно стійкі випадкові значення &
Підтримується. Для цього можна використовувати \texttt{os.urandom()}, 
а також функції та методи бібліотеки, які генерують ключі випадковим чином. 
Наприклад, \texttt{AESGCM.generate\_key()} генерує випадковий ключ необхідної довжини. &
Підтримується функцією {\footnotesize\texttt{Crypto.Random.get\_random\_bytes(n)}}, яка є криптографічно безпечною та звертається безпосередньо до системного генератора випадкових чисел ОС (\texttt{os.urandom} / \texttt{getrandom}). \\
\hline

Високо\-рівневий API &
Підтримується. Бібліотека має високорівневий API з готовими безпечними криптографічними конструкціями та низькорівневий hazmat API. 
Рекомендується використовувати high-level API, коли це можливо.&
Практично відсутній. Бібліотека орієнтована на низькорівневе та середньорівневе API, де розробник повинен самостійно контролювати параметри алгоритмів, генерувати Nonce/IV правильної довжини та вручну перевіряти теги автентифікації. \\
\hline

Апаратне прискорення &
Через OpenSSL та його апаратно оптимізовані реалізації.&
Підтримується за наявності апаратних інструкцій процесора: AES-NI та CLMUL (для прискорення обчислень GHASH у режимі AES-GCM) для архітектур x86/x64, а також ARM Crypto Extensions. За їх відсутності перемикається на оптимізовані C-процедури. \\
\hline

FIPS-сумісність / сертифікація &
Власної FIPS-сертифікації не має; 
FIPS може забезпечуватися через відповідний FIPS-валідований модуль OpenSSL.&
Власної сертифікації FIPS 140-2 / FIPS 140-3 не має. Не використовує системний OpenSSL, тому не може працювати через FIPS Object Module або FIPS Provider. \\
\hline

Ліцензія &
LICENSE.APACHE, \-LICENSE.BSD&
Подвійна ліцензія: BSD 2-Clause та Public Domain (спадщина коду вихідного PyCrypto). \\
\hline

Вразли\-вості та історія зламів &
Має історію оприлюднених вразливостей, 
які регулярно виправляються шляхом оновлення бібліотеки та її залежностей.&
Має історію виправлених вразливостей (наприклад, DoS-вразливість у реалізації KDF scrypt --- CVE-2018-15560, витоки через сайд-канали в RSA --- CVE-2023-52323). Проєкт інтегрований у Google OSS-Fuzz для превентивного виявлення дефектів пам'яті в коді на C. \\
\hline

Підтримка спільнотою &
Проєкт активно розвивається та підтримується спільнотою PyCA,
регулярно виходять оновлення й виправлення вразливостей.
Наприклад, версія 50.0.1 була випущена 25 серпня 2026 року.&
Підтримується переважно одним ключовим мейнтейнером (Helder Eijs / Legrandin) за підтримки відкритої спільноти на GitHub. Не має прямого інституційного чи корпоративного фінансування. Останній стабільний реліз 3.23.0. випущений 17 травня 2025 року, 3.24.0 наразі перебуває у розробці. \\
\hline

Актуаль\-ність документації &
Документація є актуальною та регулярно оновлюється. Доступні приклади, API-довідка та GitHub Issues. &
Документація на платформі ReadTheDocs є актуальною, містить детальний опис примітивів та готові фрагменти коду для кожного режиму роботи. Регулярно оновлюється під час нових релізів бібліотеки. \\
\hline

\end{longtable}

\subsection{Обґрунтування вибору бібліотеки}

Для виконання практичної частини та подальшої роботи обрано бібліотеку \texttt{cryptography} \cite{cryptography, cryptography git}. Основні причини вибору:

\begin{itemize}
    \item Захист від типових помилок (Misuse-resistance): наявність готових високорівневих інтерфейсів (як \texttt{AESGCM} чи \texttt{Fernet}) дозволяє шифрувати дані без ризику забути додати перевірку цілісності або помилитися з генерацією nonce. У \texttt{PyCryptodome} всі ці низькорівневі параметри доводиться контролювати вручну \cite{PyCryptodome}.
    \item Безпека пам'яті (Memory Safety): частина внутрішніх модулів бібліотеки написана на мові Rust, що суттєво знижує ризик класичних вразливостей переповнення буфера чи помилок виділення пам'яті, властивих для суто C-коду.
    \item Штатний криптографічний стандарт для Python: бібліотека підтримується офіційною організацією PyCA, постійно оновлюється, використовує системний OpenSSL і дозволяє працювати у FIPS-режимі за потреби \cite{cryptography git}.
    \item Сучасний функціонал: бібліотека активно впроваджує нові стандарти, зокрема постквантові алгоритми (ML-KEM, ML-DSA), що робить її значно перспективнішою для довгострокових проєктів порівняно з \texttt{PyCryptodome}.
\end{itemize}
